%
% graph.tex -- Template für standalone TIKZ Bilder
%
% (c) 2019 Prof Dr Andreas Müller, Hochschule Rapperswil
%
\documentclass[tikz,12pt]{standalone}
\usepackage{amsmath}
\usepackage{times}
\usepackage{txfonts}
\usepackage{pgfplots}
\usepackage{csvsimple}
\definecolor{darkgreen}{rgb}{0,0.6,0}
\definecolor{darkred}{rgb}{0.8,0,0}
\usetikzlibrary{arrows,intersections,math,calc}
\begin{document}
\def\skala{1}
\begin{tikzpicture}[>=latex,thick,scale=\skala]

\pgfmathparse{2*3.14159}
\xdef\twopi{\pgfmathresult}
\def\oben{4}
\def\unten{-4}
\fill[color=blue!10] (0,\unten) rectangle (\twopi,\oben);
\node[color=blue] at ({0.5*\twopi},0) [above] {$\Omega$};
\draw[->] (-0.1,0) -- ({2*\twopi+0.7},0) coordinate[label={$x$}];
\draw[->] (0,{\unten-0.1}) -- (0,{\oben+0.4}) coordinate[label={right:$y$}];
\begin{scope}
\clip (0,\unten) rectangle ({2*\twopi},\oben);
\foreach \y in {-4,-3.5,...,4}{
	\draw[color=darkred,line width=1.2pt] 
		plot[domain=0:90]
			({\twopi*\x/90},{\y*(sin(\x)+cos(\x))});
}
\foreach \y in {-5,-4.5,...,5}{
	\draw[color=darkred!20,line width=1.2pt] 
		plot[domain=90:180]
			({\twopi*\x/90},{\y*(sin(\x)+cos(\x))});
}
\draw (\twopi,0) -- ({2*\twopi},0);
\end{scope}
\foreach \y in {\unten,...,\oben}{
	\draw (-0.05,\y) -- (0.05,\y);
	\node at (-0.05,\y) [left] {$\y\mathstrut$};
}
\draw[color=darkgreen,line width=1.4pt] (0,\unten) -- (0,\oben);
%\node at (0,0) [left] {$0\mathstrut$};
\node at (\twopi,0) [below right] {$\displaystyle\frac{\pi}2\mathstrut$};
\draw (\twopi,-0.05) -- (\twopi,0.05);
\node at ({2*\twopi},0) [below right] {$\pi\mathstrut$};
\draw ({2*\twopi},-0.05) -- ({2*\twopi},0.05);
\draw ({1.5*\twopi},-0.05) -- ({1.5*\twopi},0.05);
\node at ({1.5*\twopi},-0.1) [below] {$\displaystyle\frac{3\pi}{4}\mathstrut$};

\end{tikzpicture}
\end{document}

